Este capítulo apresenta os conceitos que sustentam o trabalho. A Seção~\ref{sec:evolucao} discute a evolução de software e a noção de código legado. A Seção~\ref{sec:lowcode} trata das plataformas \textit{low-code} e do débito técnico associado ao Power Fx. A Seção~\ref{sec:refatoracao} aborda a refatoração e o papel dos testes unitários na alteração segura de código. Por fim, a Seção~\ref{sec:ia} discute o uso de Inteligência Artificial na Engenharia de Software, com ênfase na tradução de código e na validação de equivalência.

\section{Evolução de Software e Código Legado}
\label{sec:evolucao}

Software em uso não permanece estático. Mudanças no negócio, na legislação, na infraestrutura e nas expectativas dos usuários exigem que os sistemas sejam alterados continuamente para continuar úteis. \citeonline{sommerville2011engenharia} trata essa necessidade como parte do ciclo de vida do software e observa que, para sistemas de longa duração, os custos de evolução tendem a superar os custos de desenvolvimento inicial. O autor distingue a manutenção, que corrige defeitos e adapta o sistema a novas condições, da evolução em sentido amplo, que inclui a reestruturação do sistema para que ele continue a ser alterado com custo aceitável.

Nesse contexto surgem os sistemas legados. Para \citeonline{sommerville2011engenharia}, um sistema legado é um sistema ainda essencial para a organização, mas desenvolvido com tecnologias, estruturas ou processos que dificultam sua manutenção. Esses sistemas concentram regras de negócio que, com frequência, não estão documentadas em nenhum outro lugar: a especificação original se perdeu ou ficou desatualizada, e quem desenvolveu o sistema já não está disponível. O código passa a ser a única fonte confiável sobre o que o sistema faz.

Diante de um sistema legado, o autor descreve um espectro de estratégias: mantê-lo como está, descartá-lo, substituí-lo por um novo sistema ou submetê-lo à reengenharia. A substituição completa é a opção de maior risco, porque exige redescobrir regras que não estão registradas e reimplementá-las sem uma referência confiável de comportamento. A reengenharia, por sua vez, preserva a funcionalidade enquanto melhora a estrutura, a documentação e a manutenibilidade, e tende a apresentar menor risco e custo \cite{sommerville2011engenharia}.

\citeonline{feathers2004working} propõe uma definição mais operacional: código legado é código sem testes. O critério não é a idade nem a tecnologia, mas a ausência de um mecanismo que permita verificar, depois de uma alteração, se o comportamento existente foi mantido. Sem testes, toda modificação depende da compreensão completa do código por quem a realiza, o que raramente acontece em sistemas grandes ou mal documentados. Por isso, o autor defende que o primeiro passo para alterar um sistema legado é colocá-lo sob teste, escrevendo \textit{testes de caracterização} que registram o comportamento atual do sistema, e não o comportamento que se supõe correto. Esses testes funcionam como uma rede de proteção: qualquer alteração que mude o comportamento é detectada.

A definição de \citeonline{feathers2004working} é central para este trabalho, porque desloca o problema do legado da tecnologia para a verificabilidade. Um sistema recém-construído pode ser legado desde o início, se sua lógica não puder ser isolada e testada.

\section{Plataformas Low-Code e Débito Técnico}
\label{sec:lowcode}

Plataformas \textit{low-code} permitem construir aplicações por meio de interfaces visuais, componentes prontos e linguagens de expressão simplificadas, reduzindo a quantidade de código escrito manualmente. Sua adoção tem sido associada a ganhos de produtividade e à participação de profissionais fora das equipes de Tecnologia da Informação no desenvolvimento de soluções \cite{silva2023lowcode,silva2024lowcode}.

O Microsoft Power Apps é uma dessas plataformas. Nos aplicativos de tela (\textit{canvas apps}), o comportamento é definido em Power Fx, uma linguagem de fórmulas declarativa baseada na notação das planilhas eletrônicas \cite{powerfxdocs}. As fórmulas são escritas diretamente nas propriedades dos componentes visuais: a propriedade \texttt{OnSelect} de um botão define o que acontece ao clicar nele, a propriedade \texttt{Items} de uma galeria define quais registros são exibidos, e a propriedade \texttt{Visible} de um campo define quando ele aparece. Esse modelo reduz a barreira de entrada e acelera a construção de aplicações simples.

O mesmo modelo, porém, favorece um acoplamento forte entre regra de negócio e interface. O trecho a seguir é um exemplo ilustrativo, elaborado pela autora, de uma fórmula atribuída à propriedade \texttt{OnSelect} de um botão de inscrição em um evento:

\begin{lstlisting}[basicstyle=\ttfamily\small,frame=single,breaklines=true]
If(
    IsBlank(txtNome.Text) || dpData.SelectedDate < Today(),
    Notify("Dados invalidos", NotificationType.Error),
    If(
        CountRows(Filter(Inscricoes, EventoId = galEventos.Selected.ID))
            >= galEventos.Selected.Capacidade,
        Notify("Evento lotado", NotificationType.Warning),
        Patch(Inscricoes, Defaults(Inscricoes),
              {EventoId: galEventos.Selected.ID, Nome: txtNome.Text});
        Navigate(scrConfirmacao)
    )
)
\end{lstlisting}

Nesse trecho, três responsabilidades distintas aparecem misturadas. A regra de negócio (os dados precisam estar preenchidos, a data não pode estar no passado e o evento não pode estar lotado) depende diretamente do estado de componentes visuais, como \texttt{txtNome}, \texttt{dpData} e \texttt{galEventos}. O acesso a dados (\texttt{Filter}, \texttt{CountRows}, \texttt{Patch}) é feito na mesma expressão. A reação da interface (\texttt{Notify}, \texttt{Navigate}) encerra cada caminho da decisão. A regra não existe como uma unidade independente: ela só pode ser executada a partir da tela, com os componentes carregados e conectados à fonte de dados.

Em aplicações que crescem, esse padrão se repete em dezenas de propriedades e telas. A mesma regra pode ser reescrita em mais de um lugar, com pequenas variações, e não há um ponto único em que a lógica do sistema possa ser lida ou revisada. Como a regra depende dos componentes, não é possível testá-la isoladamente, com entradas e saídas controladas, da forma como se faz com uma função em uma linguagem de programação de propósito geral. A documentação dessas soluções também tende a ser escassa, o que já é reconhecido como desafio em equipes de desenvolvimento \textit{low-code} \cite{carlos2025lowcode}.

Esse conjunto de características pode ser entendido como débito técnico: decisões que aceleram a entrega no curto prazo, mas aumentam o custo de cada alteração futura. No caso do Power Fx, o débito não decorre de um descuido do desenvolvedor, mas do próprio modelo de programação da plataforma, que não oferece uma separação natural entre apresentação e domínio. À luz da definição de \citeonline{feathers2004working}, aplicações nesse modelo tornam-se legadas desde a sua criação, pois suas regras de negócio não podem ser colocadas sob testes unitários.

\section{Refatoração e a Importância dos Testes Unitários}
\label{sec:refatoracao}

\citeonline{fowler2018refactoring} define refatoração como uma alteração na estrutura interna do software que o torna mais fácil de entender e de modificar, sem mudar o seu comportamento observável. O ponto essencial da definição é a preservação do comportamento: refatorar não é acrescentar funcionalidades nem corrigir defeitos, mas reorganizar o código mantendo o que ele faz.

Para que essa preservação seja verificável, o autor defende duas condições. A primeira é que a refatoração seja feita em passos pequenos, cada um simples o suficiente para que seu efeito possa ser compreendido. A segunda é que exista um conjunto de testes automatizados, executado após cada passo, capaz de detectar qualquer mudança de comportamento. Quando um teste falha, o passo é desfeito ou corrigido antes de prosseguir. Essa disciplina contrasta com a reescrita de grande escala, em que muitas mudanças são feitas ao mesmo tempo e só é possível verificar o resultado no final, quando localizar a origem de um erro é muito mais difícil \cite{fowler2018refactoring}.

Os testes unitários ocupam papel central nessa abordagem. Um teste unitário verifica uma pequena parte do sistema, normalmente uma função ou classe, de forma isolada, com entradas definidas e resultados esperados explícitos. Por serem rápidos e independentes de interface, banco de dados ou rede, podem ser executados a cada alteração. \citeonline{feathers2004working} complementa essa visão ao mostrar que, em código legado, o desafio inicial é justamente criar as condições para que tais testes existam, quebrando dependências que impedem isolar a parte que se quer verificar.

Aplicado ao problema deste trabalho, isso significa extrair a regra de negócio da fórmula Power Fx para uma função que não dependa dos componentes visuais. O exemplo da Seção~\ref{sec:lowcode} poderia ser reescrito em Python como uma função que recebe valores simples e devolve uma decisão:

\begin{lstlisting}[language=Python,basicstyle=\ttfamily\small,frame=single,breaklines=true]
def validar_inscricao(nome, data_evento, hoje, inscritos, capacidade):
    if not nome or data_evento < hoje:
        return "DADOS_INVALIDOS"
    if inscritos >= capacidade:
        return "EVENTO_LOTADO"
    return "OK"
\end{lstlisting}

Separada da interface, a regra pode ser verificada diretamente por um teste unitário, por exemplo com o \textit{framework} Pytest:

\begin{lstlisting}[language=Python,basicstyle=\ttfamily\small,frame=single,breaklines=true]
def test_evento_lotado():
    resultado = validar_inscricao("Ana", date(2026, 12, 1),
                                  date(2026, 11, 1), 50, 50)
    assert resultado == "EVENTO_LOTADO"
\end{lstlisting}

A busca dos dados e as reações da interface continuam existindo, mas passam a ser responsabilidade de outras camadas, que chamam a função e tratam o seu resultado. É essa separação que permite aplicar ao sistema a refatoração incremental descrita por \citeonline{fowler2018refactoring}: cada nova alteração pode ser feita em pequenos passos, com os testes indicando se o comportamento foi preservado.

\section{Inteligência Artificial na Engenharia de Software}
\label{sec:ia}

Os Grandes Modelos de Linguagem (\textit{Large Language Models}, LLMs) passaram a ser usados em diversas tarefas da Engenharia de Software, como geração de código, explicação de trechos existentes, produção de documentação, criação de testes e tradução entre linguagens de programação. Na modernização de sistemas legados, interessa em especial a capacidade desses modelos de ler um código e descrever sua intenção em linguagem natural, o que se aproxima da engenharia reversa: recuperar, a partir da implementação, as regras que o sistema implementa.

\citeonline{pereira2026modernizacao}, ao mapear a literatura recente sobre modernização e migração automatizada de sistemas legados, indicam que abordagens apoiadas por Inteligência Artificial conseguem realizar engenharia reversa semântica do código de origem, isto é, identificar o significado das regras e não apenas a sua forma sintática. Segundo os autores, essa capacidade contribui para reduzir falhas típicas da reescrita manual e para preservar as regras de negócio durante a migração, mitigando a perda de conhecimento que acompanha as estratégias tradicionais de substituição.

A preservação, entretanto, precisa ser verificada, e não presumida. \citeonline{pan2024lost} avaliaram a tradução de código entre cinco linguagens de programação realizada por LLMs e adotaram um critério rigoroso de correção: uma tradução só foi considerada correta quando o código resultante compilava, executava e passava nos testes do programa original. Sob esse critério, o percentual de traduções corretas ficou entre 2,1\% e 47,3\%, conforme o modelo. Os autores também catalogaram os defeitos introduzidos na tradução, que incluem diferenças sutis de semântica entre as linguagens, como tratamento de tipos, de valores nulos e de bibliotecas, e mostraram que um código sintaticamente plausível pode ter comportamento diferente do original.

Os dois trabalhos são complementares para esta pesquisa. \citeonline{pereira2026modernizacao} sustentam que a IA pode recuperar e preservar regras de negócio na migração. \citeonline{pan2024lost} mostram que essa preservação só pode ser afirmada mediante validação de equivalência, ou seja, comparando o comportamento do sistema original com o do sistema migrado sob as mesmas entradas. No contexto do Power Fx, essa validação é especialmente necessária, porque a linguagem tem semântica própria para valores em branco, coerção de tipos e operações sobre tabelas, que não corresponde diretamente à do Python \cite{powerfxdocs}.

Por isso, este trabalho adota o LLM como instrumento de engenharia reversa e refatoração, mas não como garantia de correção. A equivalência entre a regra original e a implementação em Python é tratada como hipótese a ser testada, conforme os critérios descritos no Capítulo~\ref{cap:metodologia}.
